\subsubsection{Concentration of $\Gamma_t^N$}
\label{subsubsec:general-argument-bound-gamma}

We first control \(\Gamma_t^N\), the weighted moment fluctuation of the
reference empirical cloud.  The key input is that, after conditioning on the
common population-level jump randomness, the reference particles remain i.i.d.
with law \(\mu_t\).

\noindent
\textbf{A formal proof of conditionally i.i.d. property.}

We make the conditional i.i.d. structure explicit, since all concentration
bounds below are applied after conditioning on the common jump randomness.

\begin{definition}[Conditional i.i.d. reference particles]
\label{definition-conditional-iid-common-jump}
For the population field histories $\mathcal G_\ell$ fixed in the notation
section, conditional i.i.d. means
\[
\operatorname{Law}((Y_t^1,\ldots,Y_t^N)\mid\mathcal G_\ell)
=\mu_t^{\otimes N},\qquad t\in[t_\ell,t_{\ell+1}),
\]
with the pre-jump trace at $t_{\ell+1}-$. A terminal post-jump value at $T$
uses $\mathcal G_{\NofJ}$. Conditioning instead on
$\mathcal G_{\rm pop}=\mathcal G_{\NofJ}$ gives the same identity for every
time and each of its relevant traces.
\end{definition}
\begin{lemma}[The population GP jump scheme preserves conditional i.i.d.]
\label{lemma-gp-jump-implies-conditional-iid}
The population-first construction in Definition~\ref{definition-ode-and-jump}
and Assumption~\ref{assumption-common-seed} has the conditional i.i.d.
property in Definition~\ref{definition-conditional-iid-common-jump}.
\end{lemma}
\begin{proof}
Use the conditional product calculation
\eqref{eq:population-pushforward-product-law} in
\cref{lem:population-first-finite-coupling}: the common map
$\mathcal T_t$ satisfies $Y_t^i=\mathcal T_t(U_i)$ and
$\mu_t=(\mathcal T_t)_\#\mu_0$, also at each left trace.
On $[t_\ell,t_{\ell+1})$ and its final left trace, these maps and laws
are already $\mathcal G_\ell$-measurable. For bounded measurable $f_i$,
the tower property gives
\[
\begin{aligned}
 \mathbb E\left[\prod_i f_i(Y_t^i)\,\middle|\,\mathcal G_\ell\right]
 &=\mathbb E\left[\mathbb E\left[\prod_i f_i(Y_t^i)
       \,\middle|\,\mathcal G_{\rm pop}\right]\middle|\mathcal G_\ell\right]\\
 &=\mathbb E\left[\prod_i\int f_i(y)\mu_t(dy)
       \,\middle|\,\mathcal G_\ell\right]
 =\prod_i\int f_i(y)\mu_t(dy).
\end{aligned}
\]
A terminal post-jump trace uses $\mathcal G_{\NofJ}$.
The finite Gaussian coordinate is included in neither conditioning field.
\end{proof}
\paragraph{Unconditional regularity and a conditional moment event.}
In the unconditional moment assumption, write $K_0$ for the deterministic
one-particle moment bound. Conditional moment bounds will be consequences
on the event below, rather than almost-sure assumptions with the same $K_0$.
The population sigma-field $\mathcal G_{\rm pop}$ is generated by the
population fields in the population-first coupling. In particular,
$\mu_t=\operatorname{Law}(Y_t^1\mid\mathcal G_{\rm pop})$.

\begin{lemma}[Conditional moments on a population-measurable event]
\label{lem:population-conditional-moment-localization}\label{lem:population-moment-event}
Let $0<t_1<\cdots<t_{\NofJ}\le T$ be deterministic jump times. Suppose that
the paths $X^1,Y^1$ have right values and left limits at these times and
are $B$-Lipschitz on every intervening open interval, almost surely,
with the continuous branch extensions and endpoint traces of
\cref{definition-ode-and-jump},
where $B\ge0$ is deterministic. Assume that, for $S\in\{X,Y\}$,
\begin{equation}
\label{eq:population-unconditional-moment-input}
 \sup_{0\le t\le T}\sup_{q\ge1}
 q^{-1/\alpha}\bigl(\mathbb E\|S_t^1\|_{\mathbb R^d}^q\bigr)^{1/q}
 \le K_0,
 \qquad 0<\alpha\le2.
\end{equation}
The same bounds at left limits are understood by Fatou's lemma.
For $p_0\in(0,1)$, define
\[
 Q=\left\lceil(1+B)T\right\rceil+2(\NofJ+1),\qquad
 b=\max\{1,K_0\{2e(1+\alpha)\}^{1/\alpha}\},\qquad
 A_0=\frac{4Q}{p_0}.
\]
Here $b$ is a deterministic exponential-moment scale, $Q$ bounds the
number of grid time-side pairs, and $A_0$ is the threshold allocating
failure probability $p_0$ to their conditional moments.
Set $c_\alpha=2^{(\alpha-1)_+}$ and
\begin{equation}
\label{eq:population-exponential-amplitude}
 b_{\rm flow}=c_\alpha^{1/\alpha}b,\qquad
 M_{\rm pop}=A_0\exp(b^{-\alpha})\le\frac{4eQ}{p_0},
\end{equation}
and define
\begin{equation}
\label{eq:population-exponential-scale}
 k_{\rm pop}^{\alpha}
 =\frac{c_\alpha}{\log2}\{b^\alpha\log A_0+1\},
 \qquad
 K_{\rm pop}=C_\alpha k_{\rm pop},
\end{equation}
where $C_\alpha$ depends only on $\alpha$ and can be chosen so that
an exponential moment bound with scale $k$ implies the moment bound
$C_\alpha kq^{1/\alpha}$ for all $q\ge1$.
There is an event $E_{\rm pop}\in\mathcal G_{\rm pop}$ with
$\mathbb P(E_{\rm pop})\ge1-p_0$ such that, simultaneously for
$S\in\{X,Y\}$, all $t\in[0,T]$, and both sides of every jump,
\begin{equation}
\label{eq:population-fixed-scale-exponential}
 \mathbb E\!\left[
  \left.\exp\left\{\frac{\|S_t^1\|_{\mathbb R^d}^\alpha}{b_{\rm flow}^\alpha}\right\}
  \right|\mathcal G_{\rm pop}\right]\le M_{\rm pop},
\end{equation}
as well as
\begin{equation}
\label{eq:population-uniform-conditional-exponential}
 \mathbb E\!\left[
   \left.\exp\left\{\frac{\|S_t^1\|_{\mathbb R^d}^\alpha}
                              {k_{\rm pop}^\alpha}\right\}
   \right|\mathcal G_{\rm pop}\right]\le2,
\end{equation}
and hence
\begin{equation}
\label{eq:population-uniform-conditional-moments}
 \sup_{q\ge1}q^{-1/\alpha}
 \bigl(\mathbb E[\|S_t^1\|_{\mathbb R^d}^q\mid\mathcal G_{\rm pop}]\bigr)^{1/q}
 \le K_{\rm pop}.
\end{equation}
In particular $\mu_t\in\mathcal P_{\alpha,K_{\rm pop}}$ on this event.
No conditional independence of the $X$ particles is asserted or needed.
\end{lemma}

\begin{proof}
\noindent\textbf{Step 1: an unconditional exponential moment.}
For $q<1$, Lyapunov's inequality and the $q=1$ case of
\eqref{eq:population-unconditional-moment-input} give
$\mathbb E\|S_t^1\|_{\mathbb R^d}^q\le K_0^q$.
Consequently, for each integer $j\ge1$,
\[
 \mathbb E\|S_t^1\|_{\mathbb R^d}^{\alpha j}
 \le K_0^{\alpha j}\max\{1,\alpha j\}^{j}.
\]
Since $j!\ge(j/e)^j$, $\max\{1,\alpha j\}\le(1+\alpha)j$, and
$b^\alpha\ge2e(1+\alpha)K_0^\alpha$, we have
\[
 \frac{\mathbb E\|S_t^1\|_{\mathbb R^d}^{\alpha j}}{j!b^{\alpha j}}
 \le\left(\frac{eK_0^\alpha\max\{1,\alpha j\}}{jb^\alpha}\right)^j
 \le\left(\frac{eK_0^\alpha(1+\alpha)}{b^\alpha}\right)^j
 \le2^{-j}.
\]
Monotone convergence for the nonnegative exponential series therefore gives
\[
 \mathbb E\exp\{\|S_t^1\|_{\mathbb R^d}^\alpha/b^\alpha\}
 \le1+\sum_{j\ge1}2^{-j}=2.
\]
The same inequality holds at a left trace by Fatou's lemma.

\par\noindent\textbf{Step 2: one grid event and a common conditional path law.}
The deterministic branch speed $B$ is used here to transfer moment bounds
from a finite time grid to every time on each branch. It is fixed before the
random sample-localization event $E_{T,D}$ is introduced.
Set $t_0=0$ and $t_{\NofJ+1}=T$.
On each nonempty ODE branch choose a deterministic grid of mesh at most
$(1+B)^{-1}$, including the initial right value and final left trace.
The number of entries on this branch is at most
$\lceil(1+B)(t_{\ell+1}-t_\ell)\rceil+1$.
At a jump the two traces are separate entries. If $t_{\NofJ}=T$, the
zero-length final branch contributes the single post-jump value at $T$.
Thus the set $\mathcal T$ of time-side pairs satisfies
\[
 \#\mathcal T
 \le\sum_{\ell=0}^{\NofJ}
     \left(\left\lceil(1+B)(t_{\ell+1}-t_\ell)\right\rceil+1\right)
 \le(1+B)T+2(\NofJ+1)\le Q.
\]

To make the branchwise Lipschitz inequalities hold simultaneously under the
same conditional law at every grid point and time trace, use a regular
conditional law of the pair of paths $(X^1,Y^1)$ given
$\mathcal G_{\rm pop}$. Equivalently, regard each path as its finitely
many continuous branch extensions with the prescribed endpoint traces,
together with any isolated terminal post-jump value. This is a
standard Borel path space. All conditional expectations in this proof
are evaluated using this one conditional path law.
Define the countable intersection
\[
 \mathcal L:=\bigcap_{S\in\{X,Y\}}
  \bigcap_{\ell=0}^{\NofJ}
  \bigcap_{s,t\in\mathbb Q\cap(t_\ell,t_{\ell+1})}
   \{\|S_t^1-S_s^1\|_{\mathbb R^d}\le B|t-s|\}.
\]
The assumed branchwise Lipschitz property gives $\mathbb P(\mathcal L)=1$.
The tower property implies
\[
 \mathbb E[\mathbb P(\mathcal L^c\mid\mathcal G_{\rm pop})]
 =\mathbb P(\mathcal L^c)=0.
\]
Since the conditional probability is nonnegative, it is zero almost surely.
Continuity on each branch extends the rational-time inequalities to all
time pairs and to that branch's endpoint traces. Thus these inequalities
hold simultaneously under the conditional path law outside one
population-null set.

For $s\in\mathcal T$, put
\[
 Z_s^S:=\mathbb E[\exp\{\|S_s^1\|_{\mathbb R^d}^\alpha/b^\alpha\}
                         \mid\mathcal G_{\rm pop}],
 \qquad
 E_{\rm pop}:=
 \{\mathbb P(\mathcal L\mid\mathcal G_{\rm pop})=1\}
 \cap\bigcap_{S\in\{X,Y\}}\bigcap_{s\in\mathcal T}
          \{Z_s^S\le A_0\}.
\]
The event is $\mathcal G_{\rm pop}$-measurable. Step 1, the tower property,
Markov's inequality, and the finite union bound give
\[
 \mathbb P(E_{\rm pop}^c)
 \le\sum_{S\in\{X,Y\}}\sum_{s\in\mathcal T}
       \frac{\mathbb E Z_s^S}{A_0}
 \le(2Q)\frac2{A_0}=p_0.
\]

\par\noindent\textbf{Step 3: extend both exponential bounds to all times.}
Fix a population realization in $E_{\rm pop}$.
For any time-side pair $t$ on a nonempty branch, choose $s\in\mathcal T$
on that branch with $|t-s|\le(1+B)^{-1}$. Under the conditional path law,
\[
 \|S_t^1\|_{\mathbb R^d}\le\|S_s^1\|_{\mathbb R^d}+B|t-s|\le\|S_s^1\|_{\mathbb R^d}+1.
\]
For an isolated terminal value, use $s=t$.
The inequality $(u+v)^\alpha\le c_\alpha(u^\alpha+v^\alpha)$ and
$b_{\rm flow}^\alpha=c_\alpha b^\alpha$ from
\eqref{eq:population-exponential-amplitude} give
\[
 \mathbb E[\exp\{\|S_t^1\|_{\mathbb R^d}^\alpha/b_{\rm flow}^\alpha\}
                         \mid\mathcal G_{\rm pop}]
 \le e^{b^{-\alpha}}Z_s^S
 \le e^{b^{-\alpha}}A_0=M_{\rm pop}.
\]
This proves \eqref{eq:population-fixed-scale-exponential}; its fixed scale
will be retained when estimating the clipping bias.

For the normalized exponential bound, put
$\eta=\log2/\log M_{\rm pop}\in(0,1)$. By
\cref{eq:population-exponential-amplitude,eq:population-exponential-scale},
$k_{\rm pop}^\alpha=b_{\rm flow}^\alpha/\eta$.
Apply conditional Jensen directly to the all-time estimate
\eqref{eq:population-fixed-scale-exponential}:
\[
\begin{aligned}
 \mathbb E\left[e^{\|S_t^1\|_{\mathbb R^d}^\alpha/k_{\rm pop}^\alpha}
                       \,\middle|\,\mathcal G_{\rm pop}\right]
 &=\mathbb E\left[\left(e^{\|S_t^1\|_{\mathbb R^d}^\alpha/b_{\rm flow}^\alpha}\right)^\eta
                       \,\middle|\,\mathcal G_{\rm pop}\right]\\
 &\le\left(\mathbb E\left[e^{\|S_t^1\|_{\mathbb R^d}^\alpha/b_{\rm flow}^\alpha}
                       \,\middle|\,\mathcal G_{\rm pop}\right]\right)^\eta
 \le M_{\rm pop}^\eta=2.
\end{aligned}
\]
The same conditional path law is used at all time traces, proving
\eqref{eq:population-uniform-conditional-exponential} simultaneously.

\par\noindent\textbf{Step 4: obtain every ordinary conditional moment.}
Fix $q\ge1$. For $u>0$,
\[
 \frac d{du}\log\{u^qe^{-(u/k_{\rm pop})^\alpha}\}
 =\frac q u-\frac{\alpha u^{\alpha-1}}{k_{\rm pop}^\alpha}.
\]
This derivative is positive before
$u=k_{\rm pop}(q/\alpha)^{1/\alpha}$ and negative after it.
Evaluation at that maximum, also including $u=0$ by continuity, gives
\[
 u^q\le
 k_{\rm pop}^q\left(\frac q{\alpha e}\right)^{q/\alpha}
 e^{(u/k_{\rm pop})^\alpha}\qquad(u\ge0).
\]
Integrate under the same conditional path law and use
\eqref{eq:population-uniform-conditional-exponential} to obtain
\[
 \bigl(\mathbb E[\|S_t^1\|_{\mathbb R^d}^q\mid\mathcal G_{\rm pop}]\bigr)^{1/q}
 \le2^{1/q}k_{\rm pop}\left(\frac q{\alpha e}\right)^{1/\alpha}
 \le C_\alpha k_{\rm pop}q^{1/\alpha},
 \qquad
 C_\alpha=2\max\{1,(\alpha e)^{-1/\alpha}\}.
\]
This proves \eqref{eq:population-uniform-conditional-moments} for every
$q\ge1$ and every time trace on $E_{\rm pop}$.
Finally, $\mu_t=\operatorname{Law}(Y_t^1\mid\mathcal G_{\rm pop})$
implies $\mu_t\in\mathcal P_{\alpha,K_{\rm pop}}$ there.
\end{proof}

\paragraph{Probability bookkeeping.}
The event $E_{\rm pop}$ is measurable with respect to the sigma-field
under which the reference samples are conditionally iid. Thus all
conditional concentration arguments can be applied for each population
realization in this event, with the deterministic scale $K_{\rm pop}$.
If their conditional failure probability is at most $p_1$, the
unconditional failure probability, including $E_{\rm pop}^{\rm c}$, is
at most $p_0+p_1$. Conditioning on a sample-localization event is not
used to claim independence. If only the reference moment bound is
needed, the $X$ constraints can be omitted and $A_0$ replaced by $2Q/p_0$.

\noindent
\textbf{Conditional concentration of \(\Gamma_t^N\)}

We first prove a localized uniform-in-time bound: truncate the features, apply
conditional McDiarmid at grid times, and then remove the truncation.

Throughout this subsection, the jump times are deterministic, with
$0=t_0<t_1<\cdots<t_{\NofJ}\le T$ and $t_{\NofJ+1}=T$.
Every supremum and time grid includes both traces at each jump.
More precisely, the closure of an ODE interval has its post-jump value at
the left endpoint and its pre-jump value at the right endpoint.  The last
interval ends at $T$, and a jump at $T$ has an additional terminal post-jump
value.  An interval of zero length contributes that terminal value.
We condition only on the full population randomness $\mathcal G_{\rm pop}$.
Conditional on it, the reference trajectories are independent copies of
the same population trajectory and have marginal law $\mu_t$.

Let $E_{\rm pop}\in\mathcal G_{\rm pop}$ be the event supplied by
the population-moment localization lemma.  Retain both of its conclusions:
on $E_{\rm pop}$, for deterministic $K_{\rm pop}>0$, $K_0>0$,
$c_0=(K_0/b_{\rm flow})^\alpha>0$, and $M_{\rm pop}\ge1$,
\begin{equation}
\label{eq:gamma-population-localization-input}
\begin{gathered}
 \sup_{t\le T}\sup_{r\ge1}r^{-1/\alpha}
 \left(\int\|y\|_{\mathbb R^d}^r\,\mu_t(dy)\right)^{1/r}\le K_{\rm pop},
 \\
 \sup_{t\le T}\int
 \exp\!\left(c_0\left(\frac{\|y\|_{\mathbb R^d}}{K_0}\right)^\alpha\right)
 \,\mu_t(dy)\le M_{\rm pop}.
\end{gathered}
\end{equation}
Here $K_0$ is the unconditional moment scale; $K_{\rm pop}$ and
$M_{\rm pop}$ are deterministic functions of the probability budget and
the time grid.  In particular, they are not asserted to be bounded
uniformly in the realized population Gaussian fields.

Fix $D\ge1$. The deterministic weights from
\cref{eq:common-local-moment-weights} are
\[
 \gamma_m=\gamma_m(K_{\rm pop},D)=\mathrm L_Da_m.
\]
Define the radial projection onto the radius-$D$ ball by
\[
\Pi_D(y)=y\min\{1,D/\|y\|_{\mathbb R^d}\}\quad(y\ne0),\qquad \Pi_D(0)=0.
\]
As the metric projection onto a closed convex ball,
$\|\Pi_Dx-\Pi_Dy\|_{\mathbb R^d}\le\|x-y\|_{\mathbb R^d}$.
Let $\phi_m^{(D)}=\phi_m\circ\Pi_D$. The feature assumption
\cref{assumption-structured-moment-smoothness}(i),
\eqref{eq:coordinate-feature-growth}, gives
$\|\phi_m^{(D)}(y)\|_{\mathsf H_m}\le L_m(1+D^m)$.
This amplitude bounds the effect of replacing one sample in McDiarmid's
inequality. The spatial Lipschitz bound below is used separately for
time interpolation. Define the respective envelopes
\begin{equation}
\label{eq:gamma-amplitude-lipschitz-envelopes}
 A(D,K_{\rm pop})=2\sum_{m\ge1}\gamma_mL_m(1+D^m),
 \qquad
 \Lambda(D,K_{\rm pop})=
 \sum_{m\ge1}\gamma_mL_m(1+2D^{m-1}).
\end{equation}
For $\|u\|_{\mathbb R^d},\|v\|_{\mathbb R^d}\le D$, the segment $v+s(u-v)$ remains in the ball.
Integrating the first derivative in \eqref{eq:coordinate-feature-growth}
gives
\[
 \|\phi_m(u)-\phi_m(v)\|_{\mathsf H_m}
 \le\int_0^1 L_m(1+D^{m-1})\|u-v\|_{\mathbb R^d}\,ds
 =L_m(1+D^{m-1})\|u-v\|_{\mathbb R^d}.
\]
Apply this with $u=\Pi_Dx$, $v=\Pi_Dy$ and use the $1$-Lipschitz
bound above. Thus $\phi_m^{(D)}$ is globally Lipschitz with constant
at most $L_m(1+D^{m-1})\le L_m(1+2D^{m-1})$, as used in
\eqref{eq:gamma-amplitude-lipschitz-envelopes}.
Differentiability of $\Pi_D$ on the sphere is not needed.
Define also
\begin{equation}
\label{eq:gamma-conditional-tail-sum}
\begin{aligned}
 S(D,K_{\rm pop})&=
 \sum_{m\ge1}\gamma_mL_m
 \left[1+\{K_{\rm pop}(2m)^{1/\alpha}\}^{m}\right],
 \\
 \overline\beta_T(D)&=
 2S(D,K_{\rm pop})\sqrt{M_{\rm pop}}
 \exp\!\left[-\frac{c_0}{2}
               \left(\frac{D}{K_0}\right)^\alpha\right].
\end{aligned}
\end{equation}

\begin{lemma}[Localized uniform high-probability bound under the GP jump scheme]
\label{lemma-uniform-sup-gamma-bound-localized-conditional}
Assume the conditional i.i.d. conclusion of
\cref{lemma-gp-jump-implies-conditional-iid}, the feature bounds in
\cref{assumption-structured-moment-smoothness}(i), and
\cref{assumption-uniform-concentration-piecewise-lipschitz}
with deterministic constant $B$.  Assume
\eqref{eq:gamma-population-localization-input} on $E_{\rm pop}$,
and assume the sums in
\eqref{eq:gamma-amplitude-lipschitz-envelopes} and
\eqref{eq:gamma-conditional-tail-sum} are finite.
For $N\ge2$ and $p_\Gamma\in(0,1)$, set $Q_N=(\NofJ+1)(N+1)$ and
\begin{equation}
\label{eq:gamma-explicit-envelope}
 \delta_\Gamma(D)=\overline\beta_T(D)
 +\frac{A(D,K_{\rm pop})}{\sqrt N}
 +A(D,K_{\rm pop})\sqrt{\frac{\log(Q_N/p_\Gamma)}{2N}}
 +\frac{2BT\Lambda(D,K_{\rm pop})}{N}.
\end{equation}
Then
\begin{equation}
\label{eq:gamma-localized-probability}
 \mathbb P\left(
  E_{\rm pop}\cap E_{T,D}\cap
  \left\{\sup_{t\le T}\Gamma_t^N>\delta_\Gamma(D)\right\}
 \right)\le p_\Gamma.
\end{equation}
Consequently, if
$\mathbb P(E_{\rm pop}^c)\le p_{\rm pop}$ and
$\mathbb P(E_{T,D}^c)\le p_{\rm rad}$, then
\[
 \mathbb P\left(
  E_{\rm pop}\cap E_{T,D}\cap
  \left\{\sup_{t\le T}\Gamma_t^N\le\delta_\Gamma(D)\right\}
 \right)\ge1-p_{\rm pop}-p_{\rm rad}-p_\Gamma.
\]
\end{lemma}

\begin{proof}
\noindent\textbf{Step 1: concentrate the clipped empirical moments.}
Condition on $\mathcal G_{\rm pop}$, without conditioning on $E_{T,D}$,
and write $\mathbb E_{\mathcal G}[Z]=\mathbb E[Z\mid\mathcal G_{\rm pop}]$.
At a fixed time trace, let
\[
 F_t^D=\sum_{m\ge1}\gamma_m
 \left\|\frac1N\sum_{i=1}^N\phi_m^{(D)}(Y_t^i)
       -\int\phi_m^{(D)}(y)\,\mu_t(dy)\right\|_{\mathsf H_m}.
\]
If a single sample $y$ is replaced by $y'$, the reverse triangle inequality
and $\|\phi_m^{(D)}(y)\|_{\mathsf H_m}\le L_m(1+D^m)$ give
\[
 |F_t^D(y,Y^2,\ldots)-F_t^D(y',Y^2,\ldots)|
 \le\frac1N\sum_m\gamma_m
       \|\phi_m^{(D)}(y)-\phi_m^{(D)}(y')\|_{\mathsf H_m}
 \le\frac{A(D,K_{\rm pop})}{N}.
\]
For each fixed $m$, conditional independence makes every off-diagonal term
in the squared centered sample mean zero: for $i\ne j$,
\[
\begin{aligned}
 &\mathbb E_{\mathcal G}\left\langle
 \phi_m^{(D)}(Y_t^i)-\mathbb E_{\mathcal G}\phi_m^{(D)}(Y_t^i),
 \phi_m^{(D)}(Y_t^j)-\mathbb E_{\mathcal G}\phi_m^{(D)}(Y_t^j)
 \right\rangle\\
 &\qquad=\left\langle
 \mathbb E_{\mathcal G}(\phi_m^{(D)}(Y_t^i)-\mathbb E_{\mathcal G}\phi_m^{(D)}(Y_t^i)),
 \mathbb E_{\mathcal G}(\phi_m^{(D)}(Y_t^j)-\mathbb E_{\mathcal G}\phi_m^{(D)}(Y_t^j))
 \right\rangle=0.
\end{aligned}
\]
The remaining diagonal terms give
\[
\begin{aligned}
 &\mathbb E_{\mathcal G}\left\|
    \frac1N\sum_i\phi_m^{(D)}(Y_t^i)
           -\mathbb E_{\mathcal G}\phi_m^{(D)}(Y_t^1)\right\|_{\mathsf H_m}^2\\
 &\quad=\frac1N\left(
      \mathbb E_{\mathcal G}\|\phi_m^{(D)}(Y_t^1)\|_{\mathsf H_m}^2
      -\|\mathbb E_{\mathcal G}\phi_m^{(D)}(Y_t^1)\|_{\mathsf H_m}^2\right)
 \le\frac{L_m^2(1+D^m)^2}{N}.
\end{aligned}
\]
Cauchy--Schwarz bounds the expected norm by the square root of this
quantity. Summing with the nonnegative weights $\gamma_m$ gives
\[
 \mathbb E[F_t^D\mid\mathcal G_{\rm pop}]
 \le\frac1{\sqrt N}\sum_m\gamma_mL_m(1+D^m)
 \le\frac{A(D,K_{\rm pop})}{\sqrt N}.
\]
Apply \cref{lemma:mcdiarmid} with squared-difference sum
$N\{A(D,K_{\rm pop})/N\}^2=A(D,K_{\rm pop})^2/N$. It gives, for every $u>0$,
\[
 \mathbb P\left(
 F_t^D>\frac{A(D,K_{\rm pop})}{\sqrt N}+u
 \,\middle|\,\mathcal G_{\rm pop}\right)
 \le \exp\left(-\frac{2Nu^2}{A(D,K_{\rm pop})^2}\right).
\]
When $A=0$, this conclusion is read as $F_t^D=0$.
These arguments apply first to finite partial sums and then to the
infinite weighted sum by the finite envelope and monotone convergence.

\par\noindent\textbf{Step 2: extend the estimate between grid points.}
Subdivide each closed ODE interval, with its correctly labelled endpoint
traces, into $N$ equal subintervals.  The union has at most $Q_N$ nodes.
For $s,t$ on the same ODE interval, use the clipped-feature Lipschitz bound
from \eqref{eq:gamma-amplitude-lipschitz-envelopes}. The reverse triangle
inequality separates the empirical and population changes as
\[
\begin{aligned}
 |F_t^D-F_s^D|
 &\le\sum_m\gamma_mL_m(1+2D^{m-1})
 \left(\frac1N\sum_i\|Y_t^i-Y_s^i\|_{\mathbb R^d}
             +\mathbb E_{\mathcal G}\|Y_t^1-Y_s^1\|_{\mathbb R^d}\right)\\
 &\le2B\Lambda(D,K_{\rm pop})|t-s|.
\end{aligned}
\]
At each grid node substitute
$u=A(D,K_{\rm pop})\sqrt{\log(Q_N/p_\Gamma)/(2N)}$ in the preceding
concentration inequality. When $A>0$, the union over at most $Q_N$ nodes
costs
\[
 Q_N\exp\left(-\frac{2N}{A(D,K_{\rm pop})^2}
              \frac{A(D,K_{\rm pop})^2\log(Q_N/p_\Gamma)}{2N}\right)
 =p_\Gamma.
\]
When $A=0$, $F_t^D=0$ directly. Every time trace lies within $T/N$ of a
grid node on its own branch. The time Lipschitz estimate adds at most
$2BT\Lambda(D,K_{\rm pop})/N$. Consequently, for each population realization in $E_{\rm pop}$,
\begin{equation}
\label{eq:gamma-clipped-uniform-probability}
 \mathbb P\left(\sup_{t\le T}F_t^D>
 \delta_\Gamma(D)-\overline\beta_T(D)
 \,\middle|\,\mathcal G_{\rm pop}\right)\le p_\Gamma.
\end{equation}
Averaging this conditional bound preserves its failure budget.

\par\noindent\textbf{Step 3: remove clipping on the sample event.}
On $E_{T,D}$, the clipped and original empirical feature averages agree.
For $\|y\|_{\mathbb R^d}\le D$, the population integrand difference is zero; otherwise
the feature growth bound and $D^m\le\|y\|_{\mathbb R^d}^m$ give
\[
 \|\phi_m(y)-\phi_m(\Pi_Dy)\|_{\mathsf H_m}
 \le L_m(2+\|y\|_{\mathbb R^d}^m+D^m)\le2L_m(1+\|y\|_{\mathbb R^d}^m).
\]
Thus the triangle inequality gives
\begin{equation}
\label{eq:gamma-population-clipping-bias}
 \Gamma_t^N\le F_t^D+\beta_T^{\mathcal G}(D),\qquad
 \beta_T^{\mathcal G}(D)=
 2\sum_m\gamma_mL_m\sup_{s\le T}
 \int(1+\|y\|_{\mathbb R^d}^m)\mathbf1_{\{\|y\|_{\mathbb R^d}>D\}}\,\mu_s(dy).
\end{equation}
This is a conditional population integral, not the unconditional moment
of a randomly selected population field.  \Cref{lemma-beta-tail-bound}, using only the population input
\eqref{eq:gamma-population-localization-input}, gives
$\beta_T^{\mathcal G}(D)\le\overline\beta_T(D)$ on
$E_{\rm pop}$.  Intersecting the event in
\eqref{eq:gamma-clipped-uniform-probability} with $E_{\rm pop}$ and $E_{T,D}$ proves
\eqref{eq:gamma-localized-probability} and the final probability bound.
\end{proof}

\begin{lemma}[Conditional population truncation bias]
\label{lemma-beta-tail-bound}
On $E_{\rm pop}$, the bias $\beta_T^{\mathcal G}(D)$ in
\eqref{eq:gamma-population-clipping-bias} is
bounded by $\overline\beta_T(D)$ in
\eqref{eq:gamma-conditional-tail-sum}.
If only the first, moment-bound part of
\eqref{eq:gamma-population-localization-input} is retained, the valid
alternative, for $D\ge eK_{\rm pop}$, is
\begin{equation}
\label{eq:gamma-moment-only-bias}
 \beta_T^{\mathcal G}(D)\le
 2S(D,K_{\rm pop})
 \exp\!\left[-\frac12
               \left(\frac{D}{eK_{\rm pop}}\right)^\alpha\right].
\end{equation}
\end{lemma}
\begin{proof}
Write $u=c_0(D/K_0)^\alpha$. Conditional Markov, using the exponential-moment
bound \eqref{eq:population-fixed-scale-exponential} of
\cref{lem:population-moment-event}, and Cauchy--Schwarz give,
uniformly over time traces,
\[
 \mu_t(\|y\|_{\mathbb R^d}>D)\le\min\{1,M_{\rm pop}e^{-u}\},
\]
\[
\int(1+\|y\|_{\mathbb R^d}^m)\mathbf1_{\{\|y\|_{\mathbb R^d}>D\}}\,\mu_t(dy)
 \le\left[1+\{K_{\rm pop}(2m)^{1/\alpha}\}^{m}\right]
       \sqrt{M_{\rm pop}}e^{-u/2}.
\]
Indeed, Cauchy--Schwarz followed by the $L^2$ triangle inequality gives
\begin{equation}
\label{eq:gamma-tail-cauchy-schwarz}
\begin{aligned}
 \int(1+\|y\|_{\mathbb R^d}^m)\mathbf1_{\{\|y\|_{\mathbb R^d}>D\}}\,\mu_t(dy)
 &\le\left(\int(1+\|y\|_{\mathbb R^d}^m)^2\,\mu_t(dy)\right)^{1/2}
                   \mu_t(\|y\|_{\mathbb R^d}>D)^{1/2}\\
 &\le\left[1+\left(\int\|y\|_{\mathbb R^d}^{2m}\,\mu_t(dy)\right)^{1/2}\right]
                   \sqrt{\min\{1,M_{\rm pop}e^{-u}\}}.
\end{aligned}
\end{equation}
Now use \eqref{eq:gamma-population-localization-input} with $r=2m$
and $\sqrt{\min\{1,z\}}\le\sqrt z$, also when $z>1$.
Summing with weights $2\gamma_mL_m$ and recalling $S$ in
\eqref{eq:gamma-conditional-tail-sum} gives
$\beta_T^{\mathcal G}(D)\le2S(D,K_{\rm pop})\sqrt{M_{\rm pop}}e^{-u/2}
=\overline\beta_T(D)$.
For the moment-only version, set $r=(D/(eK_{\rm pop}))^\alpha\ge1$.
Markov's inequality and the moment input give
\begin{equation}
\label{eq:moment-tail-markov-substitution}
 \mu_t(\|y\|_{\mathbb R^d}>D)\le D^{-r}\int\|y\|_{\mathbb R^d}^r\mu_t(dy)
 \le\left(\frac{K_{\rm pop}r^{1/\alpha}}D\right)^r=e^{-r}.
\end{equation}
Substitute this probability into \eqref{eq:gamma-tail-cauchy-schwarz}
and sum as in \eqref{eq:gamma-population-clipping-bias}:
\[
 \beta_T^{\mathcal G}(D)
 \le2\sum_m\gamma_mL_m[1+\{K_{\rm pop}(2m)^{1/\alpha}\}^m]e^{-r/2}
 =2S(D,K_{\rm pop})e^{-r/2}.
\]
This is \eqref{eq:gamma-moment-only-bias}.
\end{proof}

\begin{lemma}[Concrete bounds for the amplitude, Lipschitz, and tail sums]
\label{lemma-concrete-bound-A-Lambda}
Let $\kappa=\kappa_{\rm dec}>\max\{1,\kappa_{\rm gr}\}$ and
$q=\kappa/(\kappa-1)$.  \Cref{lem:conditional-local-coefficient-growth,lemma-localized-moment-smoothness-consequences} give
\begin{equation}
\label{eq:gamma-coefficient-growth-explicit}
 \gamma_m(K,D)L_m\le
 C\exp\left(C\{(\log(e+D))^q+(\log(e+K))^q\}-c m^\kappa\right),
\end{equation}
with fixed $C,c>0$.  In particular,
\begin{equation}
\label{eq:gamma-three-envelope-growth}
 A(D,K)\vee\Lambda(D,K)\vee S(D,K)
 \le C\exp\left(C\{(\log(e+D))^q+(\log(e+K))^q\}\right).
\end{equation}
Here the constants do not depend on $D,K,N,T,p$.
\end{lemma}
\begin{proof}
By \cref{eq:common-local-moment-weights,eq:local-envelope-growth},
$\gamma_m(K,D)=\mathrm L_{K,D}a_m$ and
$\mathrm L_{K,D}\le Ce^{C\chi(K+D)}$.
Use \cref{lem:weighted-series-estimate},
\eqref{eq:weighted-series-summand-bound}, first without the braces
(which are at least one). Since
\[
 \log(e+K+D)\le\log(e+K)+\log(e+D),\qquad
 \chi(K+D)\le2^{q-1}\{\chi(K)+\chi(D)\},
\]
this proves \eqref{eq:gamma-coefficient-growth-explicit}, with
$c=c_a/4$ allowed. For the three sums, substitute in
\eqref{eq:weighted-series-sum-bound} as follows:
\[
\begin{aligned}
 A(D,K)&=2\mathrm L_{K,D}\sum_m a_mL_m(1+D^m)
              \le C\mathrm L_{K,D}e^{C\chi(K+D)},\\
 \Lambda(D,K)&\le2\mathrm L_{K,D}\sum_m a_mL_m(2+D^m)
              \le C\mathrm L_{K,D}e^{C\chi(K+D)},\\
 S(D,K)&=\mathrm L_{K,D}\sum_m a_mL_m
                 [1+\{2^{1/\alpha}K\}^m m^{m/\alpha}]\\
 &\le C\mathrm L_{K,D}e^{C\chi(2^{1/\alpha}K+D)}.
\end{aligned}
\]
Here $D\ge1$ gives $1+2D^{m-1}\le2(2+D^m)$.
Use \eqref{eq:weighted-series-scale-comparison} to absorb the fixed
rescaling of $K+D$, then apply the displayed bound for $\chi(K+D)$.
This proves \eqref{eq:gamma-three-envelope-growth}.
\end{proof}

\begin{lemma}[Uniform radius bound for both particle clouds]
\label{lemma:uniform-radius-bound-XY}
Suppose each one-particle marginal in both clouds obeys the unconditional
bound
$\sup_{t\le T,r\ge1}r^{-1/\alpha}
(\mathbb E\|Z_t^i\|_{\mathbb R^d}^r)^{1/r}\le K_0$, where $Z=X,Y$,
and both clouds have deterministic ODE-interval Lipschitz bound $B$.
For $D\ge2eK_0$,
\begin{equation}
\label{eq:ETD-tail}
 \mathbb P(E_{T,D}^c)\le
 2N\left(2(\NofJ+1)+\frac{2BT}{D}\right)
 \exp\left[-\left(\frac{D}{2eK_0}\right)^\alpha\right].
\end{equation}
In particular, $\mathbb P(E_{T,D}^c)\le p_{\rm rad}$ when
\begin{equation}
\label{eq:gamma-radius-choice}
 D\ge2eK_0\left[
 1+\log\frac{2N\{2(\NofJ+1)+BT/(eK_0)\}}{p_{\rm rad}}
 \right]^{1/\alpha}.
\end{equation}
\end{lemma}
\begin{proof}
Apply the Markov calculation \eqref{eq:moment-tail-markov-substitution}
of \cref{lemma-beta-tail-bound} to the unconditional law of $Z_t^i$,
with $K_{\rm pop}$ replaced by $K_0$ and $D$ by $v\ge eK_0$:
$\mathbb P(\|Z_t^i\|_{\mathbb R^d}>v)\le\exp[-(v/(eK_0))^\alpha]$.
For $B>0$, choose a grid $\mathcal T_D$ of mesh at most $D/(2B)$ on each
closed ODE branch, including both endpoint traces and any isolated
terminal post-jump value. As in \cref{lem:population-moment-event},
$\#\mathcal T_D\le2(\NofJ+1)+2BT/D$.
If $\|Z_t^i\|_{\mathbb R^d}>D$, choose a grid node $s$ on the same branch with
$|t-s|\le D/(2B)$. Then
\[
\begin{gathered}
 \|Z_s^i\|_{\mathbb R^d}\ge\|Z_t^i\|_{\mathbb R^d}-B|t-s|>D-D/2=D/2,
 \\
 E_{T,D}^c\subseteq
 \bigcup_{Z\in\{X,Y\}}\bigcup_{i=1}^N\bigcup_{s\in\mathcal T_D}
       \{\|Z_s^i\|_{\mathbb R^d}>D/2\}.
\end{gathered}
\]
For an isolated terminal value use $s=t$. Substituting $v=D/2\ge eK_0$
in the preceding tail bound and summing over this union gives
\eqref{eq:ETD-tail}; independence is not needed.
For $B=0$, one node per constant segment suffices, and the displayed
bound remains valid.  Finally, $D\ge2eK_0$ implies
$2BT/D\le BT/(eK_0)$, and
\eqref{eq:gamma-radius-choice} implies that
$(D/(2eK_0))^\alpha$ exceeds the logarithm of the union factor divided
by $p_{\rm rad}$.  This proves the final assertion.
\end{proof}

\begin{lemma}[Convenient uniform high-probability bound]
\label{lemma-uniform-sup-gamma-bound-convenient}
Assume \cref{assumption-structured-moment-smoothness,assumption-moment-growth,assumption-uniform-concentration-piecewise-lipschitz}
and the conditional product law of
\cref{lemma-gp-jump-implies-conditional-iid}. Choose the population,
radius, and concentration failure budgets each as $p/3$.
Keep the deterministic schedule size explicit and set
\[
 \Xi=\frac{3N(e+T)(\NofJ+1)(1+B)}p.
\]
Suppose the population localization lemma gives
\[
 \mathbb P(E_{\rm pop}^c)\le p/3,\qquad
 M_{\rm pop}\le C\Xi^{c_M},\qquad
 K_{\rm pop}\le C(1+\log\Xi)^{b_K},
\]
where $C,c_M,b_K$ and $K_0,c_0,\alpha$ are fixed.
For every $\varepsilon>0$, one can choose $C_D,C_\varepsilon$ depending
only on the displayed structural constants and on $\varepsilon$ such
that, with
\begin{equation}
\label{eq:def-DNTP}
 D=C_D(1+\log\Xi)^{1/\alpha},
\end{equation}
the following holds for all sufficiently large $N$:
\[
 \mathbb P\left(
  E_{\rm pop}\cap E_{T,D}\cap
  \left\{\sup_{t\le T}\Gamma_t^N(D,K_{\rm pop})
    \le C_\varepsilon\Xi^\varepsilon
          \left(N^{-1/2}+\frac{BT}{N}\right)\right\}
 \right)\ge1-p.
\]
For a uniformly bounded $B$, $BT/N$ may be replaced by $T/N$ after
changing the constant.  If $\NofJ$ is also uniformly bounded, the factor
$\Xi^\varepsilon$ is equivalent, up to constants, to
$\{N(e+T)/p\}^\varepsilon$.
\end{lemma}
\begin{proof}
\noindent\textbf{Step 1: choose the radius and control the bias.}
Increasing $C_D$ first makes \eqref{eq:gamma-radius-choice} hold with
$p_{\rm rad}=p/3$.  By \eqref{eq:gamma-three-envelope-growth},
\[
 A(D,K_{\rm pop})\vee\Lambda(D,K_{\rm pop})\vee S(D,K_{\rm pop})
 \le C\exp\{C(1+\log\log(e+\Xi))^q\}.
\]
Consequently each of these quantities is bounded by
$C_\eta\Xi^\eta$ for every fixed $\eta>0$.
The bias formula \eqref{eq:gamma-conditional-tail-sum} gives
\[
 \overline\beta_T(D)
 \le C S(D,K_{\rm pop})
 \Xi^{c_M/2}
 \exp\left[-\frac{c_0}{2}
  \left(\frac{C_D}{K_0}\right)^\alpha(1+\log\Xi)\right].
\]
Choose $C_D$ large enough that
$(c_0/2)(C_D/K_0)^\alpha\ge c_M/2+3$.
Use the subpolynomial bound with exponent $1$, so
$S(D,K_{\rm pop})\le C\Xi$. It follows that
\[
 \overline\beta_T(D)
 \le C\Xi^{1+c_M/2-(c_M/2+3)}
 =C\Xi^{-2}\le N^{-1}
\]
for sufficiently large $N$, since $\Xi\ge3eN$.

\par\noindent\textbf{Step 2: absorb the concentration logarithm and combine budgets.}
Finally, $Q_N/(p/3)\le C\Xi$ and
$\sqrt{\log(C\Xi)}\le C_\varepsilon\Xi^{\varepsilon/2}$.
Use \eqref{eq:gamma-three-envelope-growth} with exponent
$\varepsilon/2$ for $A$ and $\Lambda$. Substitution in
\cref{lemma-uniform-sup-gamma-bound-localized-conditional},
\eqref{eq:gamma-explicit-envelope}, gives
\[
\begin{aligned}
 \delta_\Gamma(D)
 &\le N^{-1}+\frac{C_\varepsilon\Xi^{\varepsilon/2}}{\sqrt N}
       (1+\Xi^{\varepsilon/2})
       +\frac{C_\varepsilon BT\Xi^{\varepsilon/2}}N\\
 &\le C_\varepsilon\Xi^\varepsilon(N^{-1/2}+BT/N).
\end{aligned}
\]
The population event, the radius bound \eqref{eq:gamma-radius-choice},
and concentration \eqref{eq:gamma-localized-probability} each cost $p/3$.
Their sum is $p$.
\end{proof}
